
\section{Dyadic 计数 \texorpdfstring{$C_1(m)$}{C\_1(m)} 的两尺度分解与误差项}

定义
$$
C_1(m)=\#\Big\{0\le N\le 2^m-1:\ c_m(N)=1\Big\}.
$$
令 $\varphi=\frac{1+\sqrt5}{2}$ 并设
$$
p=\frac{1}{\varphi^2+1}=\frac{5-\sqrt5}{10}.
$$

\subsection{主要猜想与强证据：\texorpdfstring{$C_1(m)=p2^m+O(\varphi^m)$}{C\_1(m)=p2\string^m+O(phi\string^m)}}

我们提出（并用大规模计算支持）：

\begin{conjecture}[误差尺度]
存在常数 $C>0$ 使得
$$
\abs{C_1(m)-p2^m}\le C\varphi^m,
$$
等价于
$$
\left|\frac{C_1(m)}{2^m}-p\right|=O\!\left(\left(\frac{\varphi}{2}\right)^m\right).
$$
\end{conjecture}

定义缩放误差
$$
E_m=\left(\frac{2}{\varphi}\right)^m\left(\frac{C_1(m)}{2^m}-p\right)=\frac{C_1(m)-p2^m}{\varphi^m}.
$$
在 $m=100..2000$ 的全量计算中，$E_m$ 被一个小区间夹住（强证据级事实），这与 Zeckendorf-adic odometer 的 Rokhlin tower “尾段高度 $\sim F_{m+1}$”误差控制机制一致。

\begin{table}[H]
\centering
\resizebox{\linewidth}{!}{\begin{tabular}{rlllrrll}%
\hline%
$m$&$C_1(m)$ (brief)&$b_F$ (brief)&$b_Q$ (brief)&\textbf{split}&$c_m(B_m)$&$\delta_m$&$E_m$\\%
\hline%
100&\texttt{350370008666... (30d)}&\texttt{573147844013... (21d)}&\texttt{518798640200... (21d)}&0&0&\texttt{{-}9.4825801721076295E{-}2}&\texttt{{-}8.1332094902720720E{-}2}\\%
200&\texttt{444146751870... (60d)}&\texttt{453973694165... (42d)}&\texttt{408405926068... (42d)}&0&0&\texttt{{-}1.0037534923749091E{-}1}&\texttt{{-}5.7858130100275532E{-}2}\\%
500&\texttt{904742912331... (150d)}&\texttt{225591516161... (105d)}&\texttt{106931504288... (105d)}&0&0&\texttt{{-}5.2599501032597263E{-}1}&\texttt{{-}2.0884604540555579E{-}1}\\%
1000&\texttt{296157695178... (301d)}&\texttt{703303677114... (209d)}&\texttt{305033414066... (209d)}&0&0&\texttt{{-}5.6628491504853045E{-}1}&\texttt{{-}1.5269490834711748E{-}2}\\%
1500&\texttt{969439817854... (451d)}&\texttt{219261819175... (314d)}&\texttt{273958844283... (314d)}&1&1&\texttt{2.4945987091607681E{-}1}&\texttt{{-}1.1081181125677921E{-}1}\\%
2000&\texttt{317335519468... (602d)}&\texttt{683570225957... (418d)}&\texttt{138090594785... (418d)}&0&0&\texttt{{-}7.9798623529563506E{-}1}&\texttt{9.0784730714606561E{-}2}\\%
\hline%
\end{tabular}
}
\caption{Dyadic 计数的代表性取值（表内 $C_1(m)$ 仅展示前缀与位数；完整数值在 artifacts 的 CSV 中）。}
\end{table}

\begin{table}[H]
\centering
\begin{tabular}{lll}%
\hline%
\textbf{quantity}&\textbf{min}&\textbf{max}\\%
\hline%
$\delta_m$&\texttt{{-}9.9957926526822724E{-}1}&\texttt{6.1621991627117819E{-}1}\\%
$E_m$&\texttt{{-}3.1402834882109599E{-}1}&\texttt{1.9350240316234081E{-}1}\\%
\hline%
\end{tabular}

\caption{$m\in[100,2000]$ 上 $\delta_m$ 与 $E_m$ 的范围（由脚本自动生成）。}
\end{table}

\subsection{分支线性律（经验规律）：\texorpdfstring{$E_m$}{E\_m} 与 \texorpdfstring{$\delta_m$}{delta\_m} 的近闭式关系}

设
$$
r_m=\frac{b_Q}{b_F},\qquad \delta_m=r_m-1,
$$
且 $\mathrm{split\_side}=1\iff \delta_m>0$。在该分支上观察到
$$
E_m\approx -\frac{1}{2\sqrt5}+\frac12\,\delta_m,
$$
残差常数级且显著小于主幅度。这一现象可由三值纤维的欧几里得余数结构与 Binet 公式解释：余数项给出 $\delta$ 的线性主项，取整误差给出 $\frac{1}{2\sqrt5}$ 级别的常数偏移。

